% =====================================================================
%  sections/0-head.tex — strony wstępne: tytułowa, temat po angielsku,
%  Summary, spis treści.
%
%  Układ (rozmieszczenie, odstępy, rozmiary czcionek) jest pod wymagania
%  wydziału — traktuj jak mebel na stałe. Twoja część to makra z
%  core/macros.tex: \thesisType, \topicPL, \thesisAuthor, \promotor,
%  \thesisDate i \summaryEN. Strony wstępne są nienumerowane
%  (\thispagestyle{empty}); numeracja „1" zaczyna się od Wstępu.
% =====================================================================
\begin{center}
	\fontsize{20pt}{20pt}\selectfont
	\textbf{POLITECHNIKA BIAŁOSTOCKA}
	\vspace{1.0cm}

	\fontsize{14pt}{14pt}\selectfont
	\textbf{WYDZIAŁ INFORMATYKI}
	\vspace{4.0cm}

	\fontsize{18pt}{18pt}\selectfont
	\thesisType
	\fontsize{16pt}{16pt}\selectfont

	\vspace{0.8cm}

	\fontsize{14pt}{14pt}\selectfont

	\vspace{1.6cm}

	TEMAT: \topicPL

	\vspace{1.6cm}

\end{center}
\fontsize{14pt}{14pt}\selectfont

\vspace{3,4cm}

\begin{flushright}
	\begin{tabular}{rl}
		WYKONAWCA: & \thesisAuthor \\
		           &               \\
	\end{tabular}
	\vfill
\end{flushright}

\begin{flushleft}
	\begin{tabular}{rl}
		OPIEKUN PRACY DYPLOMOWEJ: & \promotor \\
		                          &           \\
	\end{tabular}
\end{flushleft}

\vspace{2cm}

\begin{center}
	\textbf{\thesisDate}
\end{center}

\thispagestyle{empty}

\pagebreak\ \thispagestyle{empty}

\noindent Subject of diploma thesis:

\noindent \topicEN

\thispagestyle{empty}
\begin{center}
	\Large
	\textbf{Summary}

	\normalsize
	\justify
	\summaryEN
\end{center}


\pagebreak

\thispagestyle{empty}
% Karta dyplomowa od dziekanatu (PDF)? Odkomentuj — pakiet pdfpages jest
% już załadowany w preambule:
%    \includepdf[pages=2]{KartaDyplomowa.pdf}
\thispagestyle{empty}

\pagestyle{empty}
\pagebreak

\normalsize

\pagebreak
% Spis treści buduje się sam: wpisy biorą się z \section i z ręcznych
% \addcontentsline (Wstęp, Podsumowanie, Dodatki).
\addtocontents{toc}{\protect\thispagestyle{empty}}
\tableofcontents
\pagebreak
